\section{Q16 reconstruction study: full results}
\label{app:recon}

This appendix expands Section~\ref{sec:q16}. The reconstruction study compares three
predictors that differ only in the information they encode, coarse (C), reconstruction
augmented (P), and observed rich (R), against two reference baselines (a per-day uniform
prediction and a training-frequency prediction), scoring a three-class fill/adverse/no-move
target with half-Brier loss (multiclass Brier divided by two). All arms share the model
family, the training budget, and the calendar dates. Only the observed information differs.

\paragraph{Per-date losses.} Table~\ref{tab:recon} lists every predictor's half-Brier loss
for each monthly cohort, together with the two decision-relevant gaps: P$-$C (does
augmenting with reconstructed structure beat the coarse view?) and P$-$frequency (does it
beat a book-free baseline?). On the August breakout cohort ($n=691$) the augmented predictor P
scores a half-Brier loss of $0.35692$, against the coarse predictor's $0.35286$ and the
training-frequency baseline's $0.33930$. The gap to coarse is $+0.00405$ and the gap to
frequency is $+0.01762$: the augmented predictor is worse on both comparisons. The pattern is not confined to August. Across the
breakout cohorts the augmented predictor fails to beat the training-frequency baseline on
the majority of dates, and the observed-rich predictor R does no better, confirming that
the failure is about the decision target rather than about reconstruction error.

\begin{longtable}{llrrrrrrrr}
\caption{Q16 reconstruction study: per-date half-Brier loss by predictor, with the two
decision-relevant gaps. C = coarse, P = reconstruction-augmented, R = observed-rich.
Positive gaps mean the augmented predictor is worse. All values from
\texttt{experiments/results/e264\_figure\_source.json}.}
\label{tab:recon}\\
\toprule
family & date & $n$ & C & P & R & uniform & train-freq. & P$-$C & P$-$freq. \\
\midrule
\endfirsthead
\toprule
family & date & $n$ & C & P & R & uniform & train-freq. & P$-$C & P$-$freq. \\
\midrule
\endhead
\bottomrule
\endfoot
breakout & 2025-05-01 & 372 & 0.310878 & 0.309019 & 0.302571 & 0.333333 & 0.294974 & -0.001858 & 0.014046 \\
breakout & 2025-06-01 & 155 & 0.235191 & 0.221313 & 0.211069 & 0.333333 & 0.247009 & -0.013879 & -0.025697 \\
breakout & 2025-07-01 & 247 & 0.288472 & 0.286019 & 0.276986 & 0.333333 & 0.281304 & -0.002453 & 0.004715 \\
breakout & 2025-08-01 & 691 & 0.352863 & 0.356918 & 0.370796 & 0.333333 & 0.339301 & 0.004054 & 0.017616 \\
breakout & 2025-09-01 & 594 & 0.299716 & 0.300729 & 0.302344 & 0.333333 & 0.287154 & 0.001013 & 0.013575 \\
breakout & 2025-10-01 & 360 & 0.283368 & 0.286514 & 0.279983 & 0.333333 & 0.279556 & 0.003145 & 0.006958 \\
breakout & 2025-11-01 & 104 & 0.203500 & 0.199986 & 0.199041 & 0.333333 & 0.244839 & -0.003514 & -0.044853 \\
rebound & 2025-05-01 & 317 & 0.300703 & 0.293038 & 0.303007 & 0.333333 & 0.283472 & -0.007665 & 0.009566 \\
rebound & 2025-06-01 & 122 & 0.273435 & 0.272781 & 0.273255 & 0.333333 & 0.260157 & -0.000654 & 0.012624 \\
rebound & 2025-07-01 & 195 & 0.281781 & 0.285348 & 0.290173 & 0.333333 & 0.267969 & 0.003567 & 0.017379 \\
rebound & 2025-08-01 & 668 & 0.310500 & 0.310111 & 0.333572 & 0.333333 & 0.289188 & -0.000389 & 0.020923 \\
rebound & 2025-09-01 & 574 & 0.304925 & 0.304402 & 0.298784 & 0.333333 & 0.279154 & -0.000523 & 0.025248 \\
rebound & 2025-10-01 & 311 & 0.301289 & 0.300011 & 0.314831 & 0.333333 & 0.276532 & -0.001277 & 0.023480 \\
rebound & 2025-11-01 & 80 & 0.259826 & 0.252390 & 0.280369 & 0.333333 & 0.263518 & -0.007436 & -0.011127 \\
\end{longtable}


\paragraph{Decomposition.} A post-hoc decomposition separates each loss gap into a
class-mixture term (differences in the realized class frequencies across cohorts) and a
within-class term (differences in per-class loss). The earlier aggregate that had appeared
to favor rich depth is localized to the class-mixture term and to equal-versus-pooled
day weighting. Under event weighting, no learned breakout arm beats the training-frequency
reference. This is why we report the gap under a fixed weighting and against an explicit
baseline rather than as a single aggregate number: the aggregate is exactly where the
artifact hid.

\paragraph{Per-class breakdown of the headline cohort.} Table~\ref{tab:augclass} opens the
August cohort by class. The augmented predictor P and the observed-rich predictor R shift the
predicted class probabilities toward the recovered structure, but the per-class losses do not
improve over the coarse predictor C on the classes that matter for the decision. The total loss
rises accordingly. The one-hot rows confirm that the target is genuinely three-way and that no
single class dominates the loss, so the null is not an artifact of a degenerate label.

\begin{table}[h]\centering
\caption{August breakout cohort ($n=691$): per-class predicted probability ($q$) and per-class
half-Brier loss ($\ell$) for each method, with total loss. Classes are fill (F), adverse (A),
and no-move (N). Source: \texttt{experiments/results/e264\_figure\_source.json}.}
\label{tab:augclass}
\begin{tabular}{lrrrrrrr}
\toprule
method & $q_F$ & $q_A$ & $q_N$ & $\ell_F$ & $\ell_A$ & $\ell_N$ & total \\
\midrule
C & 0.367583 & 0.231548 & 0.400868 & 0.375428 & 0.680592 & 0.142871 & 0.352863 \\
P & 0.367583 & 0.231548 & 0.400868 & 0.400633 & 0.688576 & 0.125261 & 0.356918 \\
R & 0.367583 & 0.231548 & 0.400868 & 0.440539 & 0.708046 & 0.112043 & 0.370796 \\
onehot-F & 0.367583 & 0.231548 & 0.400868 & 0.000000 & 1.000000 & 1.000000 & 0.632417 \\
onehot-A & 0.367583 & 0.231548 & 0.400868 & 1.000000 & 0.000000 & 1.000000 & 0.768452 \\
onehot-N & 0.367583 & 0.231548 & 0.400868 & 1.000000 & 1.000000 & 0.000000 & 0.599132 \\
uniform & 0.367583 & 0.231548 & 0.400868 & 0.333333 & 0.333333 & 0.333333 & 0.333333 \\
train-frequency & 0.367583 & 0.231548 & 0.400868 & 0.366264 & 0.578116 & 0.176634 & 0.339301 \\
\bottomrule
\end{tabular}\end{table}


\paragraph{Controlled mechanism.} The synthetic control of Section~\ref{sec:q16} places the
same question in a setting with a known hidden state and a known optimal action. Restricted
summaries incur $0.075$ decision regret relative to the full-information optimum, but a
matched direct predictor and a posterior-restoration predictor produce equal decision
coordinates, and at $512$ training draws their outputs coincide. The lesson transfers: when
the decision does not require the hidden structure, reconstructing it spends capacity
without changing the action, and when a conditional shift is invisible in the observation
marginal it harms both approaches equally. Reconstruction is therefore neither a free lunch
nor a targeted fix. It is orthogonal to the decision unless the decision needs precisely the
recovered coordinate.

\paragraph{Finite-queue special case.} The constructed finite FIFO analysis isolates the
narrowest version of the original Q16 question: can the displayed \emph{count} of resting
orders narrow which fills are feasible beyond what aggregate volume already implies? Over
$1000$ enumerated histories, aggregate volume leaves $59$ ambiguous and adding counts leaves
$19$. Counts strictly narrow the feasible outcome in $40$ histories. The tightening is real
but rare, and on the admitted native Hyperliquid cohort the same tightening was not observed
($0$ of $961$ primary. $0$ of $962$ in a postrun completion check), consistent with the
reconstruction diagnostic: the extra information exists but seldom changes the decision.
